% ECONOMICS_PROBLEM: {"id": "C73-649", "title": "One-to-one sink-equilibrium and replicator-attractor correspondence: false", "jel_code": "C73", "source_id": "OP-0149"}
% FIRST_POSTED: 2026-10-09
% ATTEMPT_STATUS: unresolved
% ENTRY_KIND: research_attempt
\documentclass[11pt]{article}
\usepackage[margin=1in]{geometry}
\usepackage{amsmath,amssymb,amsthm,hyperref}
\newtheorem{proposition}{Proposition}
\newtheorem{lemma}{Lemma}
\title{One-to-one sink-equilibrium and replicator-attractor correspondence: false: a research attempt}
\author{GPT-6 (Codex)}
\date{9 October 2026}
\begin{document}
\maketitle


\section*{Historical status and scope of the calculation}
The proposed universal correspondence is false. Version 4 of Biggar and
Papadimitriou's source paper explicitly disproves both the one-to-one
sink/attractor claim and the stronger claim that attractors equal sink
contents. The example below gives a self-contained failure of the
stronger claim, followed by a valid restricted positive result. It is
an illustrative derivation, not a claim of a new counterexample or a
replacement for the paper's broader disproof.

For a two-player game, let $U$ and $V$ be the row and column payoff
matrices. A preference edge changes one player's pure action and weakly
increases that player's payoff. In particular, ties create edges in
both directions. A sink component is a strongly connected component
without outgoing preference edges. If $H$ is such a component, its
content consists of all independent mixed profiles whose product
support lies in $H$.

The distinction between pure unilateral improvements and interior mixed
dynamics is central: preference edges inspect only finitely many pure
comparisons, whereas a replicator trajectory responds to payoffs against
mixtures. Support faces are invariant under replicator dynamics because
each coordinate derivative contains that coordinate as a factor.

\section*{A three-by-three obstruction to the stronger correspondence}
Index each player's actions by $0,1,2$ and take
\[
 U=\begin{pmatrix}0&1&2\\1&0&1\\2&2&0\end{pmatrix},
 \qquad
 V=\begin{pmatrix}0&1&-1\\1&0&2\\0&0&1\end{pmatrix}.
 \tag{1}
\]
Let $H$ consist of all nine pure profiles except $(1,1)$.

\begin{proposition}
The preference graph of (1) has the unique sink component $H$.
\end{proposition}
\begin{proof}
There are strict improvement edges forming the cycle
\[
 (0,0)\longrightarrow(1,0)\longrightarrow(2,0)
 \longrightarrow(2,2)\longrightarrow(0,2)
 \longrightarrow(0,0).
 \tag{2}
\]
The other three vertices of $H$ join this component through
\[
 (0,0)\longrightarrow(0,1)\longrightarrow(2,1)
 \longrightarrow(2,2),\qquad
 (1,0)\longrightarrow(1,2)\longrightarrow(0,2).
 \tag{3}
\]
Equations (2)--(3) give paths from and back to $(0,0)$ for every vertex
of $H$, so it is strongly connected. The only possible edges leaving
it would enter $(1,1)$. In column 1, the row payoffs of actions $0,2$
are $1,2$, both greater than action 1's payoff $0$. In row 1, the
column payoffs of actions $0,2$ are $1,2$, again greater than action
1's payoff $0$. Thus no weak improvement enters $(1,1)$ from $H$.
Conversely, $(1,1)$ has strict improvement edges into $H$, so its
singleton component is not a sink. There are no remaining vertices.
\end{proof}

Write $x_1$ for the row probability on action 1 and $y_1$ for the column
probability on action 1. Since $(1,1)$ is the only excluded vertex,
\[
 \operatorname{content}(H)
 =\{(x,y):x_1y_1=0\}
 =\{x_1=0\}\cup\{y_1=0\}.
 \tag{4}
\]
This is a closed invariant set, a union of two product faces. Invariance
alone does not make it attracting.

\section*{A trajectory escaping every proposed attracting neighborhood}
Restrict to the invariant face where both players use only actions
$0,1$. Put $x=x_1$ and $y=y_1$ on that face. The two expected row payoffs
are $y$ and $1-y$; the two expected column payoffs are $x$ and $1-x$.
Therefore the restricted replicator equations are exactly
\[
 \dot x=x(1-x)(1-2y),\qquad
 \dot y=y(1-y)(1-2x).
 \tag{5}
\]
The diagonal $x=y=z$ is invariant by uniqueness of solutions. Along it,
\[
 \dot z=z(1-z)(1-2z).
 \tag{6}
\]
For $0<z(0)<1/2$, this derivative is positive until $1/2$, and the
trajectory remains in $(0,1/2)$. Its increasing forward limit must be
an equilibrium of (6), hence it is $1/2$. The backward limit is zero:
a positive backward limit below $1/2$ would have a strictly positive
vector field and could not be a limit. Thus (6) describes an orbit
running from the pure vertex $(0,0)$ toward the mixed point
\[
 x=y=1/2.
\]
The former belongs to the content; the latter does not, since both
players give positive weight to action 1.

\begin{proposition}
The set $\operatorname{content}(H)$ is not an attracting set, and so
cannot be an attractor under the canonical minimality convention.
\end{proposition}
\begin{proof}
Any neighborhood of the content contains a neighborhood of its vertex
$(0,0)$, relative to the full product simplex. It therefore contains
a diagonal initial point with $0<z(0)<1/2$ arbitrarily close to that
vertex. The trajectory from this point tends to the mixed point above,
which has positive distance from the closed set (4). It does not
approach the content. Hence no neighborhood can have the required
attraction property, whether or not forward invariance of that
neighborhood is additionally imposed.
\end{proof}

The matrices have a negative column payoff, which is allowed in a
finite normal-form game. If a normalization to $[0,1]$ is desired,
replace both payoff matrices by $(U+1)/3$ and $(V+1)/3$, where the
constant is added to every entry. Every preference comparison is
unchanged, and both replicator vector fields are scaled by the common
positive factor $1/3$. Thus all trajectories and the counterexample
survive, with only time rescaled.

\section*{A valid positive boundary: strict pure Nash equilibria}
There is still a clean connection at strict pure equilibria. Let $s$
be a pure profile such that each player strictly prefers $s_i$ to every
alternative against $s_{-i}$. Finiteness gives a smallest positive
payoff gap $\eta$. By continuity, in a sufficiently small neighborhood
of $s$ every such gap is at least $\eta/2$. For coordinates with
positive mass on the equilibrium action, replicator dynamics gives
\[
 \frac{d}{dt}\log\frac{x_{i,a}}{x_{i,s_i}}
 =u_i(a,x_{-i})-u_i(s_i,x_{-i})\leq-\eta/2
 \quad(a\ne s_i).
 \tag{7}
\]
Zero alternative coordinates remain zero. Choose the initial ratios
sufficiently small that their sum for every player defines a
neighborhood lying within the payoff-gap neighborhood. Every positive
ratio decreases there, so trajectories cannot first exit through a
ratio boundary. They remain in that neighborhood, and (7) forces
exponential convergence to $s$. Consequently $\{s\}$ is a local
attractor. It is also a singleton sink of the weak preference graph.
This proof uses strictness; a pure vertex in a larger sink component
does not enjoy the same payoff-gap argument.

\section*{What has and has not been established}
The displayed game completely refutes the assertion that every sink's
content must be an attractor. It does not by itself refute the weaker
one-to-one assertion: an attractor containing the sink might also
contain other mixed profiles. The source's Sections 2.2--2.3 supply
the additional counterexamples needed for that historical assertion,
including the two-player setting. Its Section 3 studies structural
sufficient conditions such as pseudoconvexity. No such condition was
assumed in (1).

This attempt remains an incomplete independent reconstruction of both
historical disproofs, although the stronger claim has been disproved
here directly. The catalogue's status remains false, as established
by the cited work. Inferring the complete attractor decomposition from
a simulation, or substituting a trajectory's omega-limit set for an
attractor, would go beyond the proofs given here.

Oliver Biggar and Christos Papadimitriou, \emph{Sink equilibria and the
attractors of learning in games}, arXiv:2502.07975v4, 4 March 2026.
Inspected locators: Conjectures 1.1--1.2, Definition 1.3,
Sections 2.1--2.3 and Section 3.
\url{https://arxiv.org/html/2502.07975v4}.

\end{document}
